%% Copernicus Publications Manuscript Preparation Template for LaTeX Submissions
%% ---------------------------------
%% This template should be used for copernicus.cls
%% The class file and some style files are bundled in the Copernicus Latex Package, which can be downloaded from the different journal webpages.
%% For further assistance please contact Copernicus Publications at: production@copernicus.org
%% https://publications.copernicus.org/for_authors/manuscript_preparation.html


%% Please use the following documentclass and journal abbreviations for preprints and final revised papers.

%% 2-column papers and preprints
\documentclass[hess, manuscript]{copernicus}


%% \usepackage commands included in the copernicus.cls:
%\usepackage[german, english]{babel}
\usepackage{tabularx}
%\usepackage{cancel}
%\usepackage{multirow}
%\usepackage{supertabular}
%\usepackage{algorithmic}
%\usepackage{algorithm}
%\usepackage{amsthm}
%\usepackage{float}
%\usepackage{subfig}
%\usepackage{rotating}
\usepackage{pdflscape}
\usepackage[section]{placeins}
\usepackage{booktabs}
\usepackage{threeparttable}
\usepackage{xcolor}
\usepackage[dvipsnames]{xcolor}
\usepackage{makecell}

\begin{document}

\title{An algorithmic framework to ascertain and compare mechanical properties from common snowpit observations}

% \Author[affil]{given_name}{surname}

\Author[1]{Mary Kate}{Connelly}
\Author[2]{Philipp}{Rosendahl}
\Author[3]{Valentin}{Adam}
\Author[1]{Samuel}{Verplanck}

\affil[1]{205 Cobleigh Hall, Montana State University, Bozeman, MT, USA 59717}
\affil[2]{TU Darmstadt}
\affil[3]{SLF Davos}

\correspondence{Mary Kate Connelly (connellymarykate@gmail.com)}

\runningtitle{Algorithmic framework for snow mechanical properties}

\runningauthor{Connelly et al.}


% ============================================================
% ABSTRACT
% ============================================================
% Key points to cover:
%   - Problem: mechanical models of snow stability require properties not routinely measured in field
%   - Approach: parameterization graph + algorithm to find/execute all calculation pathways
%   - Dataset: 10 years of SnowPilot data (50,278 pits, 371,429 layers, 14,776 ECT slabs)
%   - Key results: coverage/loss rates by method, uncertainty ranges, D11 distributions across 32 pathways
%   - Conclusion: framework is extensible; current parameterizations have high uncertainty but provide useful bounds


\section{Introduction}

% ============================================================
% 1.1 MOTIVATION
% Context: avalanche forecasting is an applied mechanics problem
% ============================================================
% - Snow stability is fundamentally a mechanical question; 70 years of mechanical models
% - Mechanical models require input parameters (mechanical properties) rarely measured directly
% - Instead, properties estimated via chains of parameterizations from field observations
% - Gap: no systematic framework to enumerate, compare, and propagate uncertainty through all possible estimation pathways
% - Motivate the SnowPilot database as a large, underutilized source of field observations
%   that could enable population-scale analysis of mechanical properties

\subsection{Mechanical Models of Snow Stability and Their Input Parameters}

% ============================================================
% 1.2 MECHANICAL MODELS AND THEIR INPUTS
% Establish what models exist and what they need
% ============================================================
% - Range of model complexity: stress-based, energy-based, combined (cite roch1965, reiweger2015,
%   heierli2008, rosendahl2020, weissgraeber2023, siron2023)
% - Focus on slab avalanche initiation: crack initiation in weak layer, propagation, slab release
% - Three categories of inputs: loading conditions (slope angle, gravity, overburden),
%   physical properties (density, layer heights), mechanical properties (E, nu, strength, toughness)
% - This paper focuses on effective isotropic elasticity (E and nu) as inputs common to
%   a wide range of current models (WEAC, etc.)
% - Introduce Hooke's law for isotropic elastic material; define E and nu
%   (equations for generalized Hooke's law and isotropic reduction already in working_draft.tex)

\subsection{Snowpits and the SnowPilot Database}

% ============================================================
% 1.3 FIELD OBSERVATIONS AND DATA SOURCE
% ============================================================
% - Define snowpit: standard field observation recording layering, grain forms,
%   densities, stability tests, terrain, weather (cite fierz2009, greene2022,
%   caatechnicalcommittee2024, norwegianwaterresourcesandenergydirectorate2022)
% - Introduce SnowPilot (cite chabot2004, chabot2016): free web application,
%   10 years of data (Sept 2014 – Sept 2025)
% - Present Table: summary statistics of the SnowPilot database
%   (total pits ~50,278, unique users ~5,381, countries ~35, layers ~371,429,
%    layers with hand hardness ~336,888, layers with grain form ~303,726,
%    layers with grain size ~176,044, layers with density ~10,468,
%    CT ~51,599, ECT ~47,684, PST ~6,213, pits near avalanche ~1,568,
%    pits on crown ~795, pits on flank ~399, pits with slope angle ~45,515)
%   NOTE: numbers from dataset_stats.ipynb; update from final dataset run

\subsection{Objectives and Scope}

% ============================================================
% 1.4 PAPER OBJECTIVES
% ============================================================
% - Present SnowPyt-MechParams: software framework implementing a parameterization
%   graph and algorithm to systematically identify and execute all valid calculation
%   pathways from snowpit observations to target mechanical parameters
% - Apply the framework to 10 years of SnowPilot data to:
%     (a) quantify data coverage (which layers/slabs have sufficient observations for each method)
%     (b) compare parameter estimates across methods
%     (c) propagate measurement uncertainty through each pathway
% - Exemplary analysis: bending stiffness D11 for ECT slabs across all 32 pathways
% - Scope: effective isotropic E and nu for slab layers; plate theory parameters
%   (A11, B11, D11, A55) for composite ECT slabs. Strength and fracture toughness
%   are out of scope but the framework is extensible to include them.
% - Briefly introduce figure showing the full chain from field observation to
%   mechanical model (intro_flowchart.pdf, bold boxes indicate scope of this paper)


\section{Data and Methods}

\subsection{Data: SnowPilot and SnowPylot}

% ============================================================
% 2.1 SNOWPILOT DATA AND PARSING
% ============================================================
% - 10 years of SnowPilot data downloaded as CAAML XML files
% - Introduce SnowPylot (cite connelly2025): open-source Python library that
%   parses CAAML XML files into Python objects (SnowPit, core_info, snow_profile,
%   stability_tests, layer, density_obs, etc.)
% - Describe the object hierarchy briefly: SnowPit → core_info (pit ID, date, user,
%   location including slope angle, country), snow_profile (layers with hardness,
%   grain form, grain size, thickness; density_obs), stability_tests (CT, ECT, PST, RB)
% - Note that only density observations aligned with stratigraphic layers are used;
%   density profile entries not matching a layer are excluded
% - Reference full SnowPylot documentation for parsing details

\subsection{Measurement Uncertainties}

% ============================================================
% 2.2 REPRESENTING FIELD OBSERVATION UNCERTAINTY
% ============================================================
% - SnowPilot stores nominal (point) values; real measurements have uncertainty
%   from both instrument/technique precision and spatial variability within the pit
% - Uncertainties are represented using the Python "uncertainties" package
%   (ufloat objects; cite lebigot2010uncertainties) to enable automatic
%   first-order uncertainty propagation through all calculations
% - Adopted uncertainties (with justification):
%     * Density (if directly observed): ±10% relative (cite conger2009, proksch2016)
%     * Slope angle: ±2° (adapted from mccammonSLOPEMEASUREMENTHUMANS2023,
%       who found ±4° for screen projections; direct slope measurement expected to be more precise)
%     * Hand hardness: ±2 sub-classes (i.e., ±0.67 HHI), e.g. a nominal 1F- spans 4F to 1F+
%       (cite holler2010; no formal recommendation given, so this is author judgement)
%     * Layer thickness: ±5% relative (no published study found; author estimate)
%     * Grain size: ±0.5 mm (author estimate)
% - Uncertainties reviewed and deemed reasonable by an experienced field observer (cite birkeland2025)
% - NOTE: "method uncertainty" (regression standard error) is tracked separately
%   from "measurement uncertainty" (input propagation); the two can be combined in
%   quadrature or reported independently. This paper reports both.

\subsection{Parameterizations for Mechanical Properties}

% ============================================================
% 2.3 PARAMETERIZATIONS
% Overview paragraph: define scope, note we select well-cited recent work,
% do not develop new parameterizations, do not include every published attempt.
% ============================================================

\subsubsection{Density}

% ============================================================
% 2.3.1 DENSITY
% Estimated from hand hardness and grain form when not directly measured
% ============================================================
% - Geldsetzer and Jamieson (2000): empirical regression from hand hardness + grain form;
%   limited to grain forms in Table (PP, PPgp, DF, RG, RGmx, FC, FCmx, DH)
%   with hand hardness ranges shown in Table tab:geldesetzer
% - Kim and Jamieson (2014) Table 2: updated regressions, larger sample size,
%   wider grain form coverage, no grain size required
% - Kim and Jamieson (2014) Table 5: alternate regression including grain size
%   as additional regressor (higher accuracy, higher data requirement)
% - Present Table: applicability matrix (grain forms × methods)
%   with hand hardness ranges (already drafted in working_draft.tex as tab:geldesetzer)
% - Note excluded grain forms: MM, SH, IF, and wet/melt forms other than MFcr
%   (no reliable parameterization exists for these forms)
% - Direct measurement (data_flow): ±10% uncertainty assumed when available

\subsubsection{Effective, Isotropic Young's Modulus}

% ============================================================
% 2.3.2 ELASTIC MODULUS
% ============================================================
% - Context: snow is anisotropic, but models such as WEAC assume effective isotropic E
% - Four parameterizations implemented (all density-based):
%     * Köchle and Schneebeli (2014): FEA of µ-CT for RG, FC, DH;
%       two exponential relationships by density range (150-250 and 250-450 kg/m³);
%       eqns already in working_draft.tex
%     * Wautier et al. (2015): FEA of X-ray CT for wide range of grain forms;
%       power law normalized by ice properties; R²=0.97; 103-544 kg/m³;
%       eqn already in working_draft.tex; note E_ice assumed 10 GPa (cf. Schottner 2026)
%     * Bergfeld et al. (2023): power law with published uncertainty in exponent (C1=4.4±0.18);
%       only parameterization with formal published uncertainty; PP, DF, RG; 110-363 kg/m³;
%       eqn already in working_draft.tex
%     * Schöttner et al. (2026): [add details from implementation]
% - Present Table: applicability matrix (grain forms × methods) with density ranges
%   (draft tab:E_applicability already in working_draft.tex; add Schöttner column)
% - Note: rewrite Köchle equations in non-dimensional form consistent with other methods

\subsubsection{Effective, Isotropic Poisson's Ratio}

% ============================================================
% 2.3.3 POISSON'S RATIO
% ============================================================
% - Two parameterizations implemented:
%     * Köchle and Schneebeli (2014): grain-form-based (no density dependence found);
%       RG: 0.171±0.026, FC: 0.130±0.040, DH: 0.087±0.063
%       (MFcr excluded; authors did not provide mean + uncertainty for that form)
%     * Srivastava et al. (2016): grain-form-based with density validity ranges;
%       RG: 0.191±0.008 (valid ρ>200 and ρ<580 kg/m³);
%       PP+DF: 0.132±0.053 (valid ρ>200 kg/m³);
%       FC+DH: 0.17±0.02 (valid ρ>200 kg/m³)
% - Note: Wautier et al. (2015) computed orthotropic Poisson's ratios but did not
%   provide an effective isotropic parameterization, so it is not implemented
% - Note: Poisson's ratio input to Poisson's ratio calculation is grain form (categorical),
%   so propagated uncertainty from input observations is zero; the stated ± values
%   reflect method (inter-sample) uncertainty only

\subsubsection{Shear Modulus}

% ============================================================
% 2.3.4 SHEAR MODULUS
% ============================================================
% - For an isotropic elastic material: G = E / (2(1+nu))
%   This is the only method implemented (derived from E and nu)
% - Wautier et al. (2015) also provides a direct shear modulus parameterization;
%   note this as an alternative / future extension
% - G is used in computing A55 (transverse shear stiffness) of the laminate

\subsubsection{Plate Theory Parameters for Composite Slabs}

% ============================================================
% 2.3.5 PLATE THEORY PARAMETERS
% Classical laminate theory following Weißgraeber and Rosendahl (2023)
% ============================================================
% - Snow slab modeled as a multi-layer plate; each layer treated as
%   homogeneous, isotropic, linearly elastic, perfectly bonded
% - Define reduced stiffness of each layer: Q_i = E_i / (1 - nu_i²)
% - A11 (extensional stiffness, N/mm): A11 = Σ Q_i * h_i
% - B11 (bending-extension coupling, N): B11 = (1/2) * Σ Q_i * (z_{i+1}² - z_i²)
% - D11 (bending stiffness, N·mm): D11 = (1/3) * Σ Q_i * (z_{i+1}³ - z_i³)
%   where z_i is the position of the bottom face of layer i measured from
%   the bottom of the slab
% - A55 (transverse shear stiffness, N/mm): A55 = κ * Σ G_i * h_i, with κ = 5/6
% - Static stress applied to weak layer by overlying slab:
%     sigma_zz = Σ rho_i * g * h_i * cos(theta)  (normal)
%     sigma_xz = Σ rho_i * g * h_i * sin(theta)  (shear)
%   (equations already in working_draft.tex)
% - Note that D11 is particularly important: stiffer, thicker slabs store
%   more elastic energy and thus drive crack propagation more effectively

\subsection{Parameterization Graph and Calculation Pathways}

% ============================================================
% 2.4 THE ALGORITHMIC FRAMEWORK
% Core methodological contribution of the paper
% ============================================================

\subsubsection{Representing Parameterizations as a Directed Graph}

% ============================================================
% 2.4.1 THE PARAMETERIZATION GRAPH
% ============================================================
% - Some parameters are inputs to other parameterizations (method chaining),
%   creating complex dependency relationships (density → E → D11)
% - Represent parameterizations as a directed acyclic graph (DAG):
%     * Parameter nodes: observed quantities (hand hardness, grain form, etc.)
%       and derived quantities (density, E, D11, etc.)
%     * Method nodes (merge nodes): combine required inputs to compute output;
%       implement AND logic (all inputs required)
%     * Edges: data flow from parameter to method to output parameter
%     * OR logic at parameter nodes: multiple incoming methods = multiple pathways
% - Present Figure: full parameterization graph (params_graph figure)
%   with node types, levels (layer vs. slab), and edge labels
% - Node levels: layer-level (density, E, nu, G) and slab-level (A11, B11, D11, A55)
% - The shared density node is a key structural feature: density feeds both
%   E-mod and Poisson's ratio (Srivastava), so the same density method must be
%   used consistently within a pathway

\subsubsection{Finding Calculation Pathways}

% ============================================================
% 2.4.2 PATHWAY ENUMERATION ALGORITHM
% ============================================================
% - Algorithm performs backward traversal from target parameter node to
%   snow_pit root node using dynamic programming (memoization)
% - At each parameter node: enumerate all incoming methods (OR logic)
% - At each merge/method node: collect all required inputs (AND logic)
% - Result: a tree of PathSegments and Branches representing one unique
%   combination of methods (a "Parameterization" / "calculation pathway")
% - Deduplication: pathways sharing identical method combinations but differing
%   in traversal structure are collapsed using a method fingerprint
% - Example: D11 has 32 unique pathways from 4 density × 4 E-mod × 2 nu methods,
%   constrained by the shared density node (density method is chosen once and
%   propagates to both E-mod and nu calculations)
% - Present figure or table listing all 32 D11 pathways with their method combinations
%   (or a representative subset to illustrate structure)

\subsubsection{Executing Parameterizations}

% ============================================================
% 2.4.3 EXECUTION ENGINE
% ============================================================
% - ExecutionEngine.execute_all(slab, target_parameter, config) runs all valid
%   pathways and returns results with success/failure for each
% - Shared intermediate results (e.g., density from the same method appearing in
%   multiple pathways) are cached; cache is cleared per slab, persists across pathways
% - This allows efficient execution: 32 D11 pathways on 14,776 slabs without
%   redundant recalculation of shared sub-paths
% - A pathway "fails" for a given layer or slab if any required input is missing
%   or falls outside a method's validity range (e.g., density out of range, grain form
%   not supported by method); failures are logged with the reason
% - Config options: include_method_uncertainty (bool) allows reporting input
%   measurement uncertainty only, or combined input + method regression uncertainty

\subsection{Constructing Slabs from ECT Observations}

% ============================================================
% 2.5 SLAB DEFINITION
% How we go from a snowpit record to a slab object for plate theory analysis
% ============================================================
% - An Extended Column Test (ECT) records failure depth and score (ECTP, ECTN, ECTV, ECTX)
% - For each ECT result, the "slab" is defined as all layers above the reported failure depth
% - If slope angle is not recorded, the pit is excluded from stress calculations
%   (but plate theory parameters A11, B11, D11 can still be computed)
% - Dataset: 14,776 ECTP slabs identified from 10 years of SnowPilot data
%   (ECTP = propagating result; ECTN, ECTV, ECTX are non-propagating)
%   NOTE: confirm whether analysis uses ECTP only or all ECT results
% - Slab-level analysis requires ALL layers in the slab to have valid inputs;
%   if any layer fails a method, the entire slab result for that pathway is NA


\section{Results}

\subsection{Coverage and Data Loss by Parameterization}

% ============================================================
% 3.1 COVERAGE AND LOSS
% How many layers/slabs can each pathway actually compute?
% ============================================================
% - Report for each method (and each pathway for multi-method parameters):
%     * Percent of layers with sufficient data to apply method ("coverage")
%     * Percent of layers where method fails ("loss"), broken down by reason
%       (missing grain form, hardness out of range, density out of range, etc.)
%     * For slab-level parameters: percent of slabs where ALL layers pass
%       (slab success rate is necessarily lower than layer success rate)
% - Present as a table and/or Sankey diagrams showing data flow from total layers
%   to successfully computed layers (Sankey figures already generated in notebooks)
% - Key numbers to report (from notebook outputs):
%     Layer coverage (density): geldsetzer 54.0%, kim_table2 63.4%, kim_table5 28.1%, measured 2.8%
%     Layer coverage (E-mod, best pathways): kim_table2→wautier 55.3%
%     Slab success (D11, best pathways): ~5.0% (e.g., geldsetzer→wautier→kochle)
%     Note the dramatic drop from layer to slab coverage (slab requires all layers to pass)

\subsection{Density Estimates Across Methods}

% ============================================================
% 3.2 DENSITY RESULTS
% ============================================================
% - Distribution of estimated density for each method applied to all SnowPilot layers
%   (violin plots or histograms, stratified by grain form)
% - Compare mean density estimates across methods: approximate range 208-283 kg/m³
% - Relative uncertainty from input measurement propagation: 10-18% depending on method
% - Note: kim_jamieson_table2 produces lower, more consistent values than geldsetzer
% - Compare estimated densities with direct measurements where available (2.8% of layers)
%   to validate parameterizations against the SnowPilot data itself
% - Figures from all_density_pathways.ipynb

\subsection{Young's Modulus Estimates Across Methods}

% ============================================================
% 3.3 ELASTIC MODULUS RESULTS
% ============================================================
% - 16 pathways (4 density × 4 E-mod methods) applied to all applicable layers
% - Distribution of E-mod estimates by pathway (violin plots, stratified by grain form)
% - Large range across methods: wautier produces lower values (~29-69 MPa),
%   kochle produces higher values (~71-252 MPa)
% - Relative uncertainty: 28-92% depending on pathway
% - Bergfeld et al. (2023) is the only method with a formally published regression
%   uncertainty; discuss the challenge of comparing uncertainties across methods
% - Figures from all_e_mod_pathways.ipynb

\subsection{Poisson's Ratio Estimates Across Methods}

% ============================================================
% 3.4 POISSON'S RATIO RESULTS
% ============================================================
% - 2 methods (kochle, srivastava), 5 pathways (because srivastava depends on density)
% - Layer coverage: kochle 42.2%, srivastava (combined) ~69.2%
% - Values: kochle ~0.144, srivastava ~0.170-0.178 depending on grain form
% - Relative input uncertainty: ~0% (grain form is categorical input; uncertainty
%   reflects method scatter, not measurement propagation)
% - Discuss interpretation: method uncertainty dominates over input uncertainty for nu
% - Figures from all_poissons_ratio_pathways.ipynb (or combine with E-mod section)

\subsection{Bending Stiffness (D11) of ECT Slabs Across All Pathways}

% ============================================================
% 3.5 D11 RESULTS - EXEMPLARY ANALYSIS
% This is the headline result: 32 pathways applied to 14,776 ECT slabs
% ============================================================
% - 32 unique pathways applied to 14,776 ECTP slabs
%   (14,776 slabs × 32 pathways = 473,032 attempted calculations)
% - Slab-level success rates: ~5% for best-covered pathways
%   (e.g., geldsetzer→wautier→kochle: 2.42×10⁸ N·m, 41.7% relative uncertainty)
% - Present Sankey diagrams showing data loss at each step of the best-covered pathway
% - Present violin plots of D11 across all successful pathways to illustrate
%   the sensitivity of D11 to method choice (3-fold range in central values)
% - Highlight the two structural sources of variability:
%     (a) density method choice strongly influences E-mod magnitude
%     (b) E-mod method choice also contributes substantially (wautier vs. kochle ~3-4x)
% - Compare D11 values stratified by ECT result (ECTP vs. ECTN/ECTV)
%   if the sample size permits statistical comparison
% - Report A11, B11, and A55 results as supplementary or in a table
% - Figures from all_D11_pathways.ipynb


\section{Discussion}

\subsection{Selecting a Preferred Calculation Pathway}

% ============================================================
% 4.1 PATHWAY SELECTION
% ============================================================
% - Framework intentionally avoids selecting a "best" pathway; instead exposes
%   the full range of outcomes so the user can make an informed choice
% - Discussion of criteria one might use to select a pathway:
%     * Coverage (which pathway computes results for the most layers/slabs?)
%     * Physical basis (which methods are validated for the grain forms present?)
%     * Uncertainty (which pathway produces the most constrained estimates?)
% - A practitioner may prefer a lower-coverage pathway with smaller uncertainty
%   over a higher-coverage pathway with large uncertainty, depending on application
% - Present recommendation: report results from multiple pathways and treat spread
%   as additional uncertainty not captured by any single method

\subsection{Uncertainty Budget: Input Measurement vs. Method Uncertainty}

% ============================================================
% 4.2 UNCERTAINTY ANALYSIS
% ============================================================
% - Two contributions to uncertainty in computed mechanical properties:
%     (a) Input measurement uncertainty (propagated via first-order Taylor expansion)
%     (b) Method (regression) uncertainty from empirical fits
% - For E-mod, relative input uncertainty alone is 28-92%; adding method uncertainty
%   makes this substantially larger
% - For nu, input uncertainty is ~0% (categorical input); method uncertainty dominates
% - For D11, combined relative uncertainty of 40-88% means order-of-magnitude differences
%   in D11 are within the plausible range even for a single pathway
% - Implication: mechanical models driven by these parameters will produce correspondingly
%   uncertain outputs; this does not invalidate the approach but it does argue against
%   using point estimates without uncertainty quantification
% - Discuss the assumption of independent layer uncertainties in slab summation
%   (spatial correlation in measurement errors could reduce or amplify total slab uncertainty)

\subsection{Limitations and Assumptions}

% ============================================================
% 4.3 LIMITATIONS
% ============================================================
% - SnowPilot data quality is heterogeneous: ~5,000 unique users with varying
%   experience levels; no formal quality-control step beyond parsing
% - Parameterizations used here cover only a subset of grain forms; layers with
%   melt forms, surface hoar, ice formations, or machine-made snow are excluded
%   from most calculations (substantial fraction of dataset in some climates)
% - Effective isotropic assumption: snow is anisotropic; mapping to isotropic
%   parameters loses information and introduces systematic bias
% - Layer thicknesses from SnowPilot are nominal; thickness variability within a pit
%   is not captured
% - Parameterizations from literature are derived from laboratory specimens;
%   representativeness for natural snowpack variability is uncertain
% - No independent validation dataset: comparing estimated vs. measured E/nu at the
%   same sites would require collocated µ-CT or acoustic measurements (not available)

\subsection{Extensions and Future Work}

% ============================================================
% 4.4 FUTURE WORK
% ============================================================
% - Framework is explicitly extensible: adding a new parameter or method requires
%   implementing the calculation function, adding a node/edge to the graph,
%   and registering the method in the dispatcher
% - Priority extensions:
%     * Weak layer failure parameters: shear strength, critical energy release rate
%       (enable full WEAC calculation from field data)
%     * Melt-freeze crust parameterizations (large fraction of dataset excluded)
%     * Anisotropic layer properties in the slab model
%     * Comparison with snow cover model outputs (SNOWPACK, CROCUS) for cross-validation
% - Potential for ML/AI: the large SnowPilot dataset with computed mechanical properties
%   could serve as training data for data-driven stability classifiers
% - Connections to Conceptual Model of Avalanche Hazard: computed D11 and static
%   stress distributions could inform problem type characterization


\section{Conclusions}

% ============================================================
% 5. CONCLUSIONS
% Numbered or bulleted; align with HESS style (prose preferred)
% ============================================================
% - The parameterization graph framework systematically enumerates all valid calculation
%   pathways from snowpit observations to target mechanical parameters, enabling
%   exhaustive method comparison at scale
% - Applied to 10 years of SnowPilot data: the framework successfully computes
%   layer-level density for 28-63% of layers and slab-level D11 for ~5% of ECT slabs,
%   depending on pathway; slab success rates are lower because all layers must pass
% - Mechanical properties can be estimated from standard field observations for
%   some grain forms and density ranges, but uncertainty is large (40-90%+ for D11);
%   this uncertainty must be communicated when these parameters feed into mechanical models
% - Current parameterizations are not sufficient to run WEAC or similar models with
%   confidence for operational avalanche forecasting; however, they can provide
%   physically-grounded bounds on slope stability when paired with stability test results
% - The framework bridges theory and practice, and is designed to grow as new
%   parameterizations become available in the literature


%% The following commands are for the statements about the availability of data sets and/or software code corresponding to the manuscript.

\codeavailability{SnowPyt-MechParams is openly available at [URL TBD]. SnowPylot is openly available at [URL TBD].}

\dataavailability{The SnowPilot database is publicly accessible at snowpilot.org. The CAAML export files used in this analysis cover the period September 1, 2014 to September 1, 2025.}

\authorcontribution{TEXT} %% this section is mandatory

\competinginterests{TEXT} %% this section is mandatory even if you declare that no competing interests are present

\begin{acknowledgements}
TEXT
\end{acknowledgements}


%% REFERENCES
\bibliographystyle{copernicus}
\bibliography{references}


\end{document}
